% !TeX root = ../thesis.tex
\chapter{Experimental Setup}

\section{How the Bellon et al. study was configured in RekenRangers}

\section{Participants and sessions}

\section{The questionnaire}
\label{sec:questionnaire}
After the final session, each participating child filled in a short paper-based questionnaire about their experience with RekenRangers. The goal was to collect structured feedback on four aspects of the tool: whether children found it motivating, whether they could use it without difficulty, whether it kept them engaged, and whether the feedback and difficulty felt appropriate.\\ \\
Designing a questionnaire for children between eight and eleven years old imposes specific constraints. Questions must be short, concrete, and free of abstract or double-barrelled formulations. Open-ended questions are impractical when the goal is comparable, structured data across a group. And the response scale must be simple enough that a child can use it without hesitation. Research on surveying children confirms that these constraints are not just practical preferences but methodological necessities: younger respondents are more sensitive to question complexity, and poorly formulated items produce unreliable data \cite{borgers2000children}.\\ \\
The questionnaire uses a 4-point Likert scale, where one star means "strongly disagree" and four stars mean "strongly agree." A 4-point scale was chosen to discourage neutral midpoint responses and encourage clearer preference indications.. The stars are presented visually on the paper form, making the scale intuitive without requiring explanation beyond a brief legend at the top of the page.\\ \\
The items are organised around four themes. Motivation is assessed through items about whether the child enjoyed using the tool and whether they would want to use it again. User friendliness is assessed through items about whether the child understood what to do and whether they needed help. Engagement is captured by asking whether RekenRangers made practising arithmetic more interesting. Feedback and difficulty alignment are assessed by asking whether the child understood the correctness feedback and whether the exercises matched their ability level.\\ \\
Two of the seven items are deliberately phrased in the negative direction: "I would not want to use RekenRangers again" and "The exercises did not match my ability well." This is a standard technique in questionnaire design to detect acquiescence bias, the tendency of respondents to agree with every statement regardless of its content. If a child gives four stars to both "I enjoyed using RekenRangers" and "I would not want to use RekenRangers again," this signals inattentive or biased responding and allows the researcher to flag that response for closer inspection.\\ \\
The items were not taken from an existing validated instrument, but are informed by established research traditions. The motivation items align with the principles measured by the Intrinsic Motivation Inventory \cite{ryan1982control}, and the user ffriendliness items reflect the core concerns of the System Usability Scale \cite{brooke1996sus}, both simplified substantially to be appropriate for young children. The full questionnaire is included in Appendix~\ref{app:questionnaire}.\\ \\